\documentclass[12pt]{article}
\usepackage{amsmath, amssymb}
\usepackage{geometry}
\geometry{margin=1in}

% Keep each labeled final-answer area together on the page.
\newcommand{\answerbox}[3][\linewidth]{%
  \begingroup
  \setlength{\fboxsep}{7pt}%
  \fbox{\begin{minipage}[t][#3][t]{\dimexpr#1-2\fboxsep-2\fboxrule\relax}
    \small\textbf{#2}\par
  \end{minipage}}%
  \endgroup
}
\newcommand{\threeanswers}[1]{%
  \par\noindent
  \answerbox[0.31\linewidth]{(a) Divergence}{0.65in}\hfill
  \answerbox[0.36\linewidth]{(b) #1}{0.65in}\hfill
  \answerbox[0.29\linewidth]{(c) Conservative?}{0.65in}\par
}
\newcommand{\curlanddivergence}{%
  \par\noindent
  \answerbox[0.56\linewidth]{Curl}{0.65in}\hfill
  \answerbox[0.40\linewidth]{Divergence}{0.65in}\par
}
\newcommand{\finalanswer}[1]{%
  \par\noindent\hfill\answerbox[0.78\linewidth]{#1}{0.65in}\par
}
\newcommand{\greenanswers}{%
  \par\noindent
  \answerbox[0.48\linewidth]{Direct line integral}{0.65in}\hfill
  \answerbox[0.48\linewidth]{Green's Theorem result}{0.65in}\par
}

\begin{document}

\noindent\textbf{MATH 261, Homework Set 7} \hfill Name: \underline{\hspace{6cm}}

\bigskip

\begin{enumerate}

% ---------------------------------------------------------
\item Let $\mathbf{F} = 2xy\,\mathbf{i} + (x^2 + y^2)\,\mathbf{j}$.
\begin{enumerate}
\item Compute the divergence of $\mathbf{F}$.
\item Compute the $k$-component of the curl of $\mathbf{F}$.
\item Is this vector field conservative?
\end{enumerate}
\vfill
\threeanswers{$k$-component of curl}
\bigskip

% ---------------------------------------------------------
\item Let $\mathbf{F} = \sin(x^2)\,\mathbf{i} + (x + e^{y^2})\,\mathbf{j}$.
\begin{enumerate}
\item Compute the divergence of $\mathbf{F}$.
\item Compute the $k$-component of the curl of $\mathbf{F}$.
\item Is this vector field conservative?
\end{enumerate}
\vfill
\threeanswers{$k$-component of curl}
\newpage

% ---------------------------------------------------------
\item Let $\mathbf{F} = x^2\,\mathbf{i} + z^2\,\mathbf{j} + xy\,\mathbf{k}$.
\begin{enumerate}
\item Compute the divergence of $\mathbf{F}$.
\item Compute the curl of $\mathbf{F}$.
\item Is this vector field conservative?
\end{enumerate}
\vfill
\threeanswers{Curl}
\bigskip

% ---------------------------------------------------------
\item Let $\mathbf{F} = x^3\,\mathbf{k}$.
\begin{enumerate}
\item Compute the divergence of $\mathbf{F}$.
\item Compute the curl of $\mathbf{F}$.
\item Is this vector field conservative?
\end{enumerate}
\vfill
\threeanswers{Curl}
\newpage

% ---------------------------------------------------------
\item Let $\mathbf{F} = yz\,\mathbf{i} + xz\,\mathbf{j} + xy\,\mathbf{k}$.
\begin{enumerate}
\item Compute the divergence of $\mathbf{F}$.
\item Compute the curl of $\mathbf{F}$.
\item Is this vector field conservative?
\end{enumerate}
\vfill
\threeanswers{Curl}
\bigskip

% ---------------------------------------------------------
\item Let $\mathbf{F}(x,y,z)
= z\cos(y)\,\mathbf{i} - xz\sin(y)\,\mathbf{j} + x\cos(y)\,\mathbf{k}$.
\begin{enumerate}
\item Compute the divergence of $\mathbf{F}$.
\item Compute the curl of $\mathbf{F}$.
\item Is this vector field conservative?
\end{enumerate}
\vfill
\threeanswers{Curl}

\newpage

% ---------------------------------------------------------
\item Find the curl and divergence for
$\mathbf{F}(x,y,z)
= x^2z\,\mathbf{i} - 2y^2z^2\,\mathbf{j} + xy^2z\,\mathbf{k}$ at the point $(1,-1,1)$.
\vfill
\curlanddivergence
\bigskip

% ---------------------------------------------------------
\item Find the curl and divergence for
$\mathbf{F}(x,y,z)
= xz^3\,\mathbf{i} - 2x^2yz\,\mathbf{j} + 2yz^4\,\mathbf{k}$ at $(1,-1,1)$.
\vfill
\curlanddivergence
\newpage

% ---------------------------------------------------------
\item Find a potential function for $\mathbf{F} = 2x\,\mathbf{i} + 3\,\mathbf{j} + 4z\,\mathbf{k}$.
\vfill
\finalanswer{Potential function}
\bigskip

% ---------------------------------------------------------
\item Find a potential function for
$\mathbf{F} = y^2\,\mathbf{i} + (2xy + e^z)\,\mathbf{j} +
ye^z\,\mathbf{k}$.
\vfill
\finalanswer{Potential function}
\newpage

% ---------------------------------------------------------
\item Find a potential function for
$\mathbf{F} = \cos z\,\mathbf{i} + 2y\,\mathbf{j} - x\sin(z)\,\mathbf{k}$.
\vfill
\finalanswer{Potential function}
\bigskip

% ---------------------------------------------------------
\item Use FTOLI to evaluate
\[
\int_C \mathbf{F} \cdot d\mathbf{r}
\]
where $\mathbf{F} = 2xyz\,\mathbf{i} + x^2z\,\mathbf{j} + x^2y\,\mathbf{k}$
and
$\mathbf{r}(t) = \langle t^2,\ \sin(\pi t/4),\ e^{t^2 - 2t} \rangle$, $0 \le t \le 2$.
\vfill
\finalanswer{Line integral}

\newpage

% ---------------------------------------------------------
\item Use FTOLI to evaluate
\[
\int_C \mathbf{F} \cdot d\mathbf{r}
\]
where $\mathbf{F} = x\,\mathbf{i} + z\cos(y)\,\mathbf{j} + \sin(y)\,\mathbf{k}$
and $C$ is the ellipse $4x^2 + 9y^2 = 36$, traversed clockwise.
\vspace{1.5in}
\finalanswer{Line integral}
\bigskip

% ---------------------------------------------------------
\item Let $\mathbf{F} = (x-y)\mathbf{i} + (x+y)\mathbf{j}$ and let $C$ be the
  circle with radius 1 traversed counterclockwise.
  Verify Green’s Theorem for $\int_C \mathbf{F} \cdot \mathbf{T}ds$
  (First compute the line integral directly then use Green's Theorem and verify
  that both methods give you the same answer)
\vfill
\greenanswers

% ---------------------------------------------------------
\newpage
\item Let $\mathbf{F} = y\mathbf{i} + x\mathbf{j}$ and let $C$ be the triangle
  with vertices $(0,0),(1,0),(1,1)$ traversed counterclockwise.
  Verify Green’s Theorem for $\int_C \mathbf{F} \cdot \mathbf{T}ds$.
\vfill
\greenanswers
\newpage


% ---------------------------------------------------------
\item Use Green’s Theorem to compute
  $\int_C \mathbf{F} \cdot \mathbf{T} ds$,
$\mathbf{F} = xy^5\,\mathbf{i} + (y^2 + x^2)\mathbf{j}$
on the rectangle with vertices $(0,0),(0,1),(1,0),(1,1)$, traversed
counterclockwise.
\vfill
\finalanswer{Circulation}
\bigskip

% ---------------------------------------------------------
\item Now compute
  $\int_C \mathbf{F} \cdot \mathbf{N} ds$
for the same $\mathbf{F}$ and rectangle of the previous exercise.
\vfill
\finalanswer{Flux}
\newpage

% ---------------------------------------------------------
\item Use Green’s Theorem to compute
  $\int_C \mathbf{F} \cdot \mathbf{T} ds$,
$\mathbf{F} = ye^x\,\mathbf{i} + (y^2 + x^2)\mathbf{j}$
on the rectangle with vertices $(0,0),(0,3),(5,0),(5,3)$, traversed
counterclockwise.
\vfill
\finalanswer{Circulation}
\bigskip

% ---------------------------------------------------------
\item Repeat the previous exercise but now compute
  $\int_C \mathbf{F} \cdot \mathbf{N} ds$ instead.
\vfill
\finalanswer{Flux}
\newpage

% ---------------------------------------------------------
\item Use Green’s Theorem to evaluate
  $\int_C \mathbf{F} \cdot \mathbf{T} ds$,
$\mathbf{F} = (x^2 + y)\mathbf{i} + (y^2 - x)\mathbf{j}$
on the triangle $(0,0),(-1,0),(-1,1)$, traversed counterclockwise.
\vfill
\finalanswer{Circulation}
\bigskip

% ---------------------------------------------------------
\item Repeat the previous exercise but now compute
  $\int_C \mathbf{F} \cdot \mathbf{N} ds$ instead.
\vfill
\finalanswer{Flux}
\newpage

% ---------------------------------------------------------
\item Use Green’s Theorem to find the work done by
$\mathbf{F}(x,y) = (x+y)\mathbf{i} + xy^2\mathbf{j}$
in moving a particle from the origin along the x-axis to (1,0), then along a
line segment to (0,1), and then back to the origin along the y-axis.
\vfill
\finalanswer{Work}
\bigskip
% ---------------------------------------------------------
\item Use Green’s Theorem to find the work done by
$\mathbf{F}(x,y) = x^2 y\,\mathbf{i} + xy^2\mathbf{j}$
in moving a particle down from the origin along the line $y=-x$ to the point
(1,-1), then up along a line segment to (1,1), and then back to the origin along
the line $y=x$.
\vfill
\finalanswer{Work}
\newpage

% ---------------------------------------------------------
\item Use Green's Theorem to find the flux of the vector field $\mathbf{F}(x,y)
  = x^2\mathbf{i} + y^2\mathbf{j}$ across the circle $x^2 + y^2 = 1$.
\vfill
\finalanswer{Flux}
\bigskip

% ---------------------------------------------------------
\item Use Green's Theorem to find the flux of
  $\mathbf{F}(x,y) = x\,\mathbf{i} + y^2\,\mathbf{j}$ across the curve obtained
  by following the graph of $y=x^2$ from (0,0) to (1,1), the line $y=1$ from
  (1,1) to (0,1), and the line $x=0$ from (0,1) to (0,0).
\vfill
\finalanswer{Flux}
\end{enumerate}

\end{document}
